\section{A discrete Gauss--Vandermonde product}
\label{dhr:sec:determinant}

For $r\geq1$ and $0\leq n_1<\cdots<n_r$, put
\[
 \Delta(\bm n)=\prod_{1\leq i<j\leq r}(n_j-n_i),\qquad
 w_n(a,c;b)=\frac{\Gamma(n+a)\Gamma(n+c)}
 {\Gamma(n+b)\Gamma(n+1)}.
\]
Define
\begin{equation}\label{dhr:eq:partition}
 Z_r(a,c;b)=\sum_{0\leq n_1<\cdots<n_r}
       \Delta(\bm n)^2\prod_{i=1}^r w_{n_i}(a,c;b),
 \qquad Z_0=1.
\end{equation}
The ordering is part of the definition. There is no additional $r!$
in any product formula below.

\begin{theorem}[Gauss--Vandermonde product]\label{dhr:thm:gauss}
Let $\Rea a,\Rea c>0$ and $\Rea d>2r-2$, where $d=b-a-c$.
Then \eqref{dhr:eq:partition} converges absolutely and locally normally,
and
\begin{equation}\label{dhr:eq:gauss-product}
 \boxed{\displaystyle
 Z_r(a,c;b)=\prod_{j=0}^{r-1}
 \frac{j!\,\Gamma(a+j)\Gamma(c+j)\Gamma(d-r+1-j)}
 {\Gamma(b-a-j)\Gamma(b-c-j)}.}
\end{equation}
For positive real parameters in this domain all summands are positive.
\end{theorem}

We give the determinant argument in detail, both to make the normalizing
factor checkable and to avoid relying on an unproved discrete-Selberg
analogy. Gamma ratio estimates and Gauss summation are classical
\cite{dlmf:gamma,dlmf:gauss}; the determinant technique belongs to the
standard evaluation toolkit \cite{krattenthaler}.

\begin{lemma}[A rational factorial determinant]\label{dhr:lem:rational-det}
For indeterminates $x_0,\ldots,x_{r-1},t$,
\begin{equation}\label{dhr:eq:rational-det}
 \det_{0\leq i,j<r}\frac{(x_i)_j}{(x_i+t)_j}
 =\frac{\displaystyle\prod_{j=0}^{r-1}(t)_j
                   \prod_{0\leq i<k<r}(x_k-x_i)}
        {\displaystyle\prod_{i=0}^{r-1}(x_i+t)_{r-1}}.
\end{equation}
The identity holds at every specialization where the displayed
rational functions are defined, and elsewhere by cancellation.
\end{lemma}
\begin{proof}
Multiply row $i$ by $(x_i+t)_{r-1}$. The resulting column polynomials
are
\[
 p_j(x)=(x)_j(x+t+j)_{r-1-j},\qquad 0\leq j<r.
\]
They have degree at most $r-1$, so their evaluation determinant is
the Vandermonde in the $x_i$ times their coefficient determinant.
To compute that coefficient determinant, evaluate at $x_i=-i$.
The evaluation matrix is lower triangular, since $(-i)_j=0$ for
$j>i$, with diagonal
\[
 p_i(-i)=(-1)^i i!\,(t)_{r-1-i}.
\]
Its determinant divided by
$\prod_{i<k}((-k)-(-i))=(-1)^{r(r-1)/2}\prod_{k=0}^{r-1}k!$
is $\prod_{j=0}^{r-1}(t)_j$. Dividing the row factors proves
\eqref{dhr:eq:rational-det} as a rational identity.
\end{proof}

\begin{proof}[Proof of Theorem~\ref{dhr:thm:gauss}]
On compact subsets of the stated domain,
$w_n=O((n+1)^{-1-\Rea d})$. For ordered nonnegative indices,
\[
 |\Delta(\bm n)|^2\leq\prod_{j=1}^r(n_j+1)^{2(j-1)}.
\]
Dropping the ordering gives a product of convergent majorant series
when $\Rea d>2r-2$. This proves absolute local normal convergence,
including that of moments of degrees at most $2r-2$.

Apply finite Cauchy--Binet to the two monic polynomial bases
$\fall{n}{i}$ and $(n+a)_j$, then pass to the normally convergent
limit. Both evaluation determinants equal the Vandermonde, so
\begin{equation}\label{dhr:eq:moment-det}
 Z_r=\det_{0\leq i,j<r}M_{ij},\qquad
 M_{ij}=\sum_{n\geq0}\fall{n}{i}(n+a)_j w_n.
\end{equation}
Shifting $n=i+k$ and applying the ordinary Gauss sum gives
\begin{equation}\label{dhr:eq:mixed-moment}
 M_{ij}=\frac{\Gamma(a+i+j)\Gamma(c+i)\Gamma(d-i-j)}
 {\Gamma(b-a-j)\Gamma(b-c)}.
\end{equation}
Indeed the hypergeometric numerator parameters are $a+i+j,c+i$,
the denominator parameter is $b+i$, and the excess is $d-i-j$.
Its real part is positive throughout the present index range.

Factor from row $i$ the quantity
\[
 R_i=\frac{\Gamma(a+i)\Gamma(c+i)\Gamma(d-i)}
 {\Gamma(b-a)\Gamma(b-c)},
\]
and from column $j$ the factor $(1-b+a)_j$. The remaining entry is
$(a+i)_j/(1-d+i)_j$. Lemma~\ref{dhr:lem:rational-det} applies with
$x_i=a+i$ and $t=1-b+c$. Since $x_k-x_i=k-i$, its Vandermonde
is $\prod_{j=0}^{r-1}j!$. The factorial identities
\[
 (1-b+a)_j=(-1)^j\frac{\Gamma(b-a)}{\Gamma(b-a-j)},\quad
 (1-b+c)_j=(-1)^j\frac{\Gamma(b-c)}{\Gamma(b-c-j)}
\]
and
\[
 (1-d+i)_{r-1}=(-1)^{r-1}
       \frac{\Gamma(d-i)}{\Gamma(d-i-r+1)}
\]
produce exactly \eqref{dhr:eq:gauss-product}; the total sign is one.
Perform the calculation first away from the auxiliary determinant's
poles, then extend by the identity theorem. The original sum and the
final product are holomorphic on the asserted domain.
\end{proof}

\begin{remark}[The one-variable boundary case]
For $r=1$, \eqref{dhr:eq:gauss-product} is simply
\[
 \sum_{n\geq0}w_n(a,c;b)
 =\frac{\Gamma(a)\Gamma(c)\Gamma(b-a-c)}
 {\Gamma(b-a)\Gamma(b-c)}.
\]
The subsequent $r=1$ resonance is the familiar zero-balanced Gauss
case. The continuation of interest is the uniform mechanism for all $r$,
not a claim of novelty for this boundary case.
\end{remark}
